\section{Reproduction and Compile-Time Contract}
\label{app:reproduction}

\subsection{Canonical command graph}

The release verifier authenticates the committed source and historical
snapshot bytes. Fresh measurements can be generated from the repository root with
\path{tools/plan_fa4_measurements.py}. The planner selects the historical
snapshot under \path{reproduction/snapshots/forward_cfc06dad} for the
noncausal study and the root causal source for the training study. It checks
supplied causal artifacts and external assets, then emits a fresh build and run
graph; it does not download data, run CUDA, or submit a job. Newly built
extensions are checked against generated bundle manifests or recorded by
content digest when the generated command executes.

List the complete graph with:
\begin{lstlisting}[language=bash]
python3 tools/plan_fa4_measurements.py list
\end{lstlisting}
For example, this validates and prints the noncausal forward graph:
\begin{lstlisting}[language=bash]
python3 tools/plan_fa4_measurements.py print \
  --family noncausal-forward \
  --python /absolute/path/python3 \
  --output-root /absolute/path/new-results \
  --noncausal-build-root /absolute/path/new-build \
  --cuda-home /absolute/path/cuda-13.0 \
  --cutlass-dsl-root /absolute/path/cutlass-dsl
\end{lstlisting}
The output and build roots must be new or empty. A blocked node is printed as a
comment and makes the planner exit with status 2; missing evidence is never
replaced by a guessed command. The generated noncausal suite preallocates
outputs, cycles provider order, and records every timing window.

\subsection{Fixed-input model evaluations}

ViT, BERT masked-language-model, and SST-2 replacement measurements use the
\code{downstream} family. Each generated process runs exactly one task with an
explicit extension root and one authenticated model/dataset pair:
\begin{lstlisting}[language=bash]
python3 tools/plan_fa4_measurements.py print \
  --family downstream \
  --python /absolute/path/python3 \
  --output-root /absolute/path/new-results \
  --noncausal-build-root /absolute/path/new-build \
  --cuda-home /absolute/path/cuda-13.0 \
  --cutlass-dsl-root /absolute/path/cutlass-dsl \
  --external-assets-manifest /absolute/path/assets.json
\end{lstlisting}
The asset manifest binds immutable revisions and every local file by byte count
and SHA256. These commands produce new, fully pinned measurements. The
historical paper rows remain receipt-backed because their original external
asset revisions were not recorded.

The \code{vit-mae} and \code{wan} families use the same asset contract. The
ViT-MAE manifest must identify the recorded 100 COCO validation images. Wan
keeps the logical model identifier separate from the authenticated local model
snapshot and builds the fast and accurate policy bundles before evaluation. The
paired HAO-BF16/TK comparison is available; HAO low-precision Wan controls remain
blocked because their exact extension build identity was not preserved.

The committed Wan table is a historical artifact. Some calibration and
provider JSON inputs needed to regenerate it are absent, so the release does
not advertise \path{build_tables.py} as a complete rebuild. The retained
receipts support the reported values, while a fresh \code{wan} run produces
replacement evidence under the current authenticated protocol.

\subsection{Offline artifacts and external inputs}

Tables, figures, and manuscript outputs supported by committed inputs are
rebuilt without network access using:
\begin{lstlisting}[language=bash]
python3 tools/reproduce_fa4_paper.py --run --offline all
\end{lstlisting}
\path{docs/fa4_measurement_reproduction.md} defines the external-asset and
artifact-manifest schemas. It also records which historical measurements are
receipt-only and which can be replaced by a fully authenticated fresh run.

\subsection{B300 aggregate reconstruction}

The source release contains the committed aggregate B300 summary used by the
paper, but not the larger raw cluster-capture archive. Regenerate the LaTeX
tables from that committed summary with:
\begin{lstlisting}[language=bash]
python3 results/fp4_fa4_b300_tuning_20260802/\
build_summary.py --from-summary
\end{lstlisting}
This route performs no network access and validates the summary schema before
rendering. Recomputing the aggregate from raw captures remains blocked until a
checksummed raw-capture bundle is published separately. Once available,
restore it beneath the B300 result directory and run
\code{build\_summary.py} without \code{--from-summary}. Cluster submission
files, scheduler identifiers, private storage locations, and secret wiring are
operational metadata rather than scientific inputs and are excluded from the
public reproduction path.

\subsection{Required operand contract}

Both retained policies require folded K64 Q/K scales selected by the MSE
policy:
\begin{lstlisting}[language=bash]
--nv-qk-fold-k64-scales both \
--nv-qk-fold-scale-select mse
\end{lstlisting}
\code{Accurate} also requires the same fixed 32-row permutation for K and V.
Pairing a folded-scale binary with ordinary block-16 Q/K scales can produce
plausible timing and invalid accuracy; the suite couples the binary and
operand arguments.

\subsection{Retained compile-time policy}

\code{Fast} uses a 12-stage K/V ring, all-affine mode 23, four-word
represented denominators, paired scale reuse, wide three-input maxima,
Q-scale preloading, early asynchronous P-scale handoff, and a 200/208/56
register split. \code{Accurate} uses 13 K/V stages, native EX2 samples, a
32-row anchor, and correction-warpgroup normalization.

Historical NV/NV, intermediate NV/MX, sparse, sampled-denominator,
direct-code, quadratic, two-CTA, alternate-owner, and interleaving probes
remain compile-gated and default-off.

\section{Terminology}

\begin{description}
  \item[Direct-P] The retained forward-inference method: NVFP4 Q/K,
  MXFP4 P/V, direct log-score-to-E2M1 probability conversion, and
  normalization from the represented probability consumed by PV.
  \item[Quantized causal backward] The retained training method: reconstruct
  P from the saved quantized Q/K payload, scales, and LSE, then use
  range-appropriate, matrix-oriented FP8 operands for gradient products.
  \item[BF16 / FP8 / FP4] Bfloat16 and 8-/4-bit floating-point families.
  E$x$M$y$ denotes $x$ exponent bits and $y$ explicit fraction bits.
  \item[SM / CTA] A streaming multiprocessor is one GPU compute unit; a
  cooperative thread array is a CUDA thread block scheduled on an SM.
  \item[TMEM / TMA] Tensor memory is Blackwell's on-chip matrix-accumulator
  scratchpad.  The Tensor Memory Accelerator moves tiles between GPU memory
  spaces.
  \item[MMA] A tensor-core matrix multiply--accumulate operation.
  \item[FMA / SFU / EX2] Fused multiply--add; special-function unit; and the
  hardware base-two exponential instruction.
  \item[GQA / RoPE / LSE] Grouped-query attention; rotary positional
  embedding; and the per-row log-sum-exp softmax normalizer.
  \item[Quarter] One N32 fragment read from an N128 score tile. Four producer
  quarters feed two K64 scaled-FP4 PV commands.
  \item[Score bank] A 128-column FP32 TMEM region receiving QK and later
  hosting transient P/scale overlays.
  \item[Output bank] A permanent 128-column FP32 PV accumulator for one query
  stage.
  \item[Represented denominator] A normalizer summed from the E2M1 payload
  and decoded block scale actually consumed by PV.
  \item[Early publication] Signaling a legal first K64 payload and scale pair
  before tail P construction completes.
  \item[Speed-of-light diagnostic] An intentionally incorrect or
  semantically altered kernel that removes work to bound latency.
  \item[Full FP4 attention] Q, K, P, and V inside the attention kernel use
  E2M1 payloads with block scales. NVFP4 QK plus FP8 PV is a control, not full
  FP4 attention.  The term makes no claim about learned QKV/O projections or
  an end-to-end training recipe.
\end{description}
